\DocumentMetadata{
  pdfversion=2.0,
  pdfstandard=ua-2,
  lang=en-US,
  tagging=on,
  tagging-setup={math/setup=mathml-SE,tabsorder=structure}
}
\documentclass[11pt,letterpaper]{article}
\usepackage[margin=0.68in,includefoot]{geometry}
\usepackage{newcomputermodern}
\usepackage{amsmath,amscd,mathtools}
\usepackage{tikz,adjustbox}
\providecommand{\square}{\mdlgwhtsquare}
\usepackage[hidelinks]{hyperref}
\hypersetup{
  pdftitle={Algebra First-Year Exam, May 2021, Part 2},
  pdfauthor={University of Florida Department of Mathematics},
  pdfsubject={Graduate mathematics examination},
  pdfdisplaydoctitle=true,
  bookmarksopen=true
}
\setcounter{secnumdepth}{0}
\setlength{\parindent}{0pt}
\setlength{\parskip}{0.28em}
\setlength{\emergencystretch}{3em}
\setlength{\leftmargini}{2.15em}
\setlength{\leftmarginii}{2.65em}
\setlength{\leftmarginiii}{3em}
\pagestyle{empty}
\raggedbottom
\allowdisplaybreaks
\NewTaggingSocketPlug{block/list/label}{examproblem}{%
  \tagstructbegin{tag=Lbl}\tagstructend%
  \tagstructbegin{tag=\LBody}%
  \tagstructbegin{tag=\ProblemHeadingTag,title-o={\ProblemHeadingText}}%
  \tagmcbegin{tag=\ProblemHeadingTag,actualtext={\ProblemHeadingText}}%
  #2%
  \tagmcend%
  \tagstructend%
}
\newenvironment{examproblems}[2]{%
  \def\ProblemHeadingTag{#2}%
  \AssignTaggingSocketPlug{block/list/label}{examproblem}%
  \begin{enumerate}%
  \setcounter{enumi}{#1}%
  \renewcommand{\labelenumi}{\textbf{\theenumi.}}%
  \setlength{\topsep}{0.35em}%
  \setlength{\partopsep}{0pt}%
  \setlength{\itemsep}{0.48em}%
  \setlength{\parsep}{0.18em}%
}{\end{enumerate}}
\newenvironment{examsubparts}{%
  \AssignTaggingSocketPlug{block/list/label}{default}%
  \begin{enumerate}%
  \setlength{\topsep}{0.18em}%
  \setlength{\partopsep}{0pt}%
  \setlength{\itemsep}{0.16em}%
  \setlength{\parsep}{0.1em}%
}{\end{enumerate}}
\newenvironment{examinstructions}{%
  \par\begingroup\itshape
}{%
  \par\endgroup\vspace{0.18em}
}
\makeatletter
\renewcommand\subsection{\@startsection{subsection}{2}{0pt}%
  {-1.25ex plus -.25ex minus -.1ex}%
  {0.55ex plus .1ex}%
  {\normalfont\large\bfseries}}
\renewcommand{\@maketitle}{%
  \newpage\null\vskip -1em%
  {\footnotesize\bfseries
    \href{https://www.ufl.edu/}{University of Florida}, \href{https://math.ufl.edu/}{Department of Mathematics}\par}%
  \vskip -0.5em%
  \tagstructbegin{tag=H1,title={Algebra First-Year Exam, May 2021, Part 2}}%
  \tagpdfparaOff%
  \tagmcbegin{tag=H1}%
    {\LARGE\bfseries\href{https://gma.math.ufl.edu/exams/algebra-first-year/}{Algebra First-Year Exam}, May 2021, Part 2\par}%
  \tagmcend%
  \tagpdfparaOn%
  \tagstructend%
  \vskip 0.25em%
  \tagmcbegin{artifact}%
    \hrule height 1pt%
  \tagmcend%
  \vskip 1em%
}
\makeatother
\title{Algebra First-Year Exam, May 2021, Part 2}
\author{}
\date{}
\begin{document}
\maketitle
\thispagestyle{empty}
\begin{examinstructions}
Answer four problems. Write your solutions clearly in complete English sentences. You may quote results (within reason) as long as you state them clearly.
\end{examinstructions}
\begin{examproblems}{0}{H2}
\def\ProblemHeadingText{Problem 1.}
\item
Prove the Chinese remainder theorem: Let~\(R\) be a commutative ring with~\(1\) and let \(I,J\) be ideals in~\(R\) such that \(I+J=R\). Then \(IJ=I\cap J\) and there is a ring isomorphism \(R/IJ\cong(R/I)\times(R/J)\).
\def\ProblemHeadingText{Problem 2.}
\item
Prove that the ring \(\mathbb{Z}[\sqrt{-2}]=\{a+b\sqrt{-2}:a,b\in\mathbb{Z}\}\) is a Euclidean domain with respect to the field norm
\[
N:\mathbb{Z}[\sqrt{-2}]\to\mathbb{Z}_{\geq 0},\qquad N(a+b\sqrt{-2})=a^2+2b^2.
\]
\def\ProblemHeadingText{Problem 3.}
\item
Prove that each of the following is irreducible in the indicated ring.
\begin{examsubparts}
\item[\textnormal{(a)}]
\(2X^3+2X^2+3X+1\in\mathbb{Q}[X]\)
\item[\textnormal{(b)}]
\(Y^3+X^2+XY-Y\in\mathbb{Q}[X,Y]\)
\item[\textnormal{(c)}]
\(X^4-X^3+4X^2-15\in\mathbb{Z}[X]\)
\end{examsubparts}
\def\ProblemHeadingText{Problem 4.}
\item
Let~\(F\) be a field and let~\(V\) be a vector space over~\(F\). Let \(S\subset V\) be a set which spans~\(V\). Use Zorn’s lemma to prove that~\(V\) has a basis which is contained in~\(S\).
\def\ProblemHeadingText{Problem 5.}
\item
Determine the number of similarity classes of matrices \(A\in\mathrm{M}_5(\mathbb{C})\) satisfying \(A^3=0\). Give a representative for each similarity class in Jordan canonical form.
\def\ProblemHeadingText{Problem 6.}
\item
This problem has two parts.
\begin{examsubparts}
\item[\textnormal{(a)}]
Let \(L/K\) be a field extension of degree~\(5\) and let \(f(X)\in K[X]\) have degree~\(4\). Prove that if \(\alpha\in L\) is a root of \(f(X)\), then \(\alpha\in K\).
\item[\textnormal{(b)}]
Give an example of a field extension \(L/K\) of degree~\(4\), a polynomial \(f(X)\in K[X]\) of degree~\(5\), and a root~\(\alpha\) of \(f(X)\) such that \(\alpha\in L\) but \(\alpha\notin K\).
\end{examsubparts}
\end{examproblems}
\end{document}
